% ============================================================================
%  constitution_v1_2_proposed.tex
%  PROPOSAL — NOT AUTHORITATIVE.
%
%  Finding F11 instrument: clause-level addressability of the Constitution.
%  This file renders Ledger Framework Constitution v1.1 with an identifier
%  C-<sec>.<n> on EVERY normative clause. It is pure addressing: not one word
%  of Constitution prose is changed. The three PARKED amendments (F2/F10/F12
%  = PARK-2/3/4) are folded in as clearly marked v1.2 deltas — the only prose
%  deltas; everything else is verbatim v1.1 + labels.
%
%  Authority is one-way and does NOT turn here. The Constitution v1.1
%  (LedgerManifesto/leddeger_manifesto.md) remains the sole authority; this
%  proposal is cited by nothing in the specification except the parking index.
%  Compile: latexmk -pdf constitution_v1_2_proposed.tex
% ============================================================================
\documentclass[11pt,a4paper]{article}

\usepackage[utf8]{inputenc}
\usepackage[T1]{fontenc}
\usepackage{lmodern}
\usepackage[margin=1in]{geometry}
\usepackage{parskip}
\usepackage{fancyhdr}
\usepackage{enumitem}

\setlength{\emergencystretch}{3em}

% --- running header: unmissable, on every page ---
\pagestyle{fancy}
\fancyhf{}
\fancyhead[C]{\small\bfseries PROPOSED --- NOT AUTHORITATIVE}
\fancyfoot[C]{\small\thepage}
\renewcommand{\headrulewidth}{0.4pt}

% --- a normative clause: bold address, then verbatim v1.1 prose ---
\newcommand{\clause}[2]{\par\medskip\noindent\textbf{#1}\hspace{0.6em}#2}

% --- a parked-amendment delta box ---
\newcommand{\pk}[3]{%
  \par\smallskip\noindent
  \fbox{\parbox{\dimexpr\linewidth-2\fboxsep-2\fboxrule\relax}{%
    \small\textbf{[PARKED AMENDMENT --- #1]}\par\smallskip
    \emph{Clause replaced (verbatim, Constitution v1.1):} #2\par\smallskip
    \emph{Proposed v1.2 replacement (verbatim):} #3}}\par\smallskip}

\begin{document}

\begin{center}
\fbox{\parbox{0.95\linewidth}{\centering
  \bfseries PROPOSED --- NOT AUTHORITATIVE.\\
  Cited by nothing in the specification except the parking index.\\[3pt]
  \normalfont\small
  This document renders Ledger Framework Constitution v1.1 with a clause-level
  address \textbf{C-\textit{sec}.\textit{n}} on every normative clause --- pure
  addressing, \emph{not one word of Constitution prose changed} --- and folds in
  the three parked amendments (F2, F10, F12) as marked v1.2 deltas. The
  Constitution v1.1 (\texttt{LedgerManifesto/leddeger\_manifesto.md}) remains the
  sole authority; authority is one-way and does not turn here.}}
\end{center}

\begin{center}
\large\bfseries The Ledger Framework --- Architectural Constitution\\
\normalsize\normalfont Clause-Addressed Rendering (proposed v1.2) of v1.1
\end{center}

\noindent\emph{Non-normative material omitted from clause numbering:} the
Abstract and the Amendment Record of v1.1 are summary and historical prose
respectively; they carry no requirement, definition, principle, or guarantee of
their own and so bear no clause identifier. Every normative clause of
\S1--\S14, Scope, and Authority is reproduced below, verbatim, with its address.

% ===================================================================
\section*{\S1 --- The Objective and the Problem}
% ===================================================================

\clause{C-1.1}{The objective of the Ledger project is to build a clean post-trade
and risk-management platform: one system of record from which positions,
valuations, profit and loss, collateral, and reports are all derived --- exactly,
deterministically, and reproducibly.}

\clause{C-1.2}{The root cause is architectural: \textbf{the same fact is stored
more than once.}}

\clause{C-1.3}{Two stores of one fact maintained by independent logic will
eventually disagree, and no amount of process removes a failure mode the
architecture itself provides.}

\clause{C-1.4}{There is exactly one canonical record --- an ordered, append-only
stream of atomic transactions --- and every other number is computed on demand as
a projection of that stream: balances, valuations, profit-and-loss statements,
balance sheets, regulatory reports.}

\clause{C-1.5}{Two projections of one record cannot disagree about any fact the
record contains.}

% ===================================================================
\section*{\S2 --- The Six Commitments}
% ===================================================================

\clause{C-2.1}{\textbf{Full auditability.} Every number is traceable, mechanically
and completely, to the recorded events that produced it. History is never
rewritten: a mistake is repaired by a recorded correction that names what it
supersedes, so the error and its repair both survive.}

\clause{C-2.2}{\textbf{Full reproducibility.} Given the same record, any party ---
the firm, a counterparty, an auditor, a regulator --- recomputes the same state,
the same views, the same contract outputs, to the last minor unit.}

\clause{C-2.3}{\textbf{Time travel embedded in the design.} The exact state of the
system at any past moment is reconstructible from the record alone, at any time,
forever --- independently of whether the software that produced the events still
exists.}

\clause{C-2.4}{\textbf{Correctness.} Illegal states are made unrepresentable
wherever possible: a check can be forgotten or bypassed; a type cannot. Where
wrongness cannot be prevented --- a contract can always compute the wrong
economics --- it is made mechanically detectable and deterministically
repairable. Defaults fail closed: the system stops rather than improvises.}

\clause{C-2.5}{\textbf{Testability.} The ledger is correct because it is built to
be tested. Every step is a pure, versioned function with explicit inputs, explicit
outputs, and strong types --- typing alone eliminates whole classes of defects
before a single test runs --- and every core invariant is expressed as an
executable check over generated products, events, and histories, never defended
only in prose. Correctness then composes: because the steps are pure and their
composition is fixed --- the map, then the fold --- verifying each part verifies
the whole. A property that cannot be tested mechanically is redesigned until it
can be. This is not an afterthought: the architecture is deliberately shaped so
that its core guarantees are executable --- property-based testing is a
construction principle, not a phase at the end.}

\clause{C-2.6}{\textbf{Clarity.} One name per component and one component per name;
behaviour carried as declared data rather than hidden in code; every claim stated
so that it can be checked by a reader with no prior exposure.}

% ===================================================================
\section*{\S3 --- The Whole Architecture in One Picture}
% ===================================================================

\clause{C-3.1}{Everything this framework does can be stated as a map followed by a
fold. A \emph{map} processes a list element by element: given a way to turn an a
into a b, it carries [a] to [b], one element at a time. A \emph{fold} accumulates
a list into a single result: each element of [a] is processed in turn and absorbed
into a running value, [a] to b.}

\clause{C-3.2}{A \textbf{unit} is anything that can be held or owed --- a stock, a
bond, an option, a currency; a \textbf{wallet} is where units are held, in signed
quantities called balances; a \textbf{transaction} is the atomic bundle of changes
the ledger records; and a \textbf{projection} is any view computed from the
record.}

\clause{C-3.3}{Every unit declares, from its terms, the conditions its life
requires watching, and the Event Monitor evaluates those watches against the clock
and the recorded observations, emitting an event onto the record the moment a
condition is met.}

\clause{C-3.4}{The smart contracts are the mapping functions: each is a stateless
piece of code that, given one event and the current state of the ledger, computes
the transaction the event requires. The Events Executor applies the map over the
list, E = [e] to T = [tx]: for each event it fires the responsible contract at the
right time, hands it the event together with the current projection of the ledger,
and collects the proposed transaction. Many contracts, one Events Executor; the
contracts carry all the product knowledge, the Events Executor carries none.}

\clause{C-3.5}{The Transaction Executor performs the fold, T = [tx] to Ledger: it
serialises, validates, and applies each transaction, and the immutable log is
precisely the list being folded. Every view --- a balance, a valuation, a report
--- is a further fold of the same list.}

\clause{C-3.6}{Three kinds of thing appear in this picture, and they must never be
confused: data, functions, and machines. The data are Event, Transaction, and
Ledger. The functions are the arrows between the data. The machines are what
evaluate the arrows --- one machine per arrow.}

\clause{C-3.7}{The clock is the single input that is not on the record, and the
Event Monitor is the single place it is consulted. The event the Monitor emits is
recorded, so everything downstream of it is a function of the record alone ---
which is exactly why a late firing produces the identical transaction. Because
watches are declared and the Monitor is deterministic, even the emission of events
is verifiable: given the record and the clock, anyone can recompute which events
should have been emitted, and when.}

\clause{C-3.8}{Both contract and apply can refuse --- a contract may find nothing
due; apply rejects a malformed transaction and writes nothing --- and a refusal is
a returned value, never a half-applied state.}

\clause{C-3.9}{Map and fold therefore advance together, one event at a time, in a
single total order.}

\clause{C-3.10}{The map is proposed while the fold is authoritative: a transaction
joins the list only when the Transaction Executor admits it, and economic
verification is simply re-running the map on the recorded event and comparing ---
which is why both the events and the transactions are on the record.}

\clause{C-3.11}{Replay re-reads transactions and never re-executes contracts, so
history cannot be silently reinterpreted; a wrong interpretation is repaired by a
compensating transaction, in the open. Verification lives with the events;
authority lives with the transactions.}

% ===================================================================
\section*{\S4 --- The Objects}
% ===================================================================

\clause{C-4.1}{At the highest level, a ledger is three things and a record: units
--- what can be held; wallets --- where it is held; balances --- how much of each
unit each wallet holds; and the immutable log that records every change to them.}

\clause{C-4.2}{A \textbf{unit} is anything that can be held or owed inside the
ledger: a currency (USD in cents), an equity security (identified by its ISIN), a
bond, a listed or OTC derivative, a securities loan, a managed-account mandate, a
margin obligation under a collateral agreement (a CSA), or any other contractual
instrument. Units are first-class and extensible: adding a new instrument type
adds a new element to the universe of units, never a modification of the
machinery. Because every unit --- asset, liability, or contractual obligation ---
is represented uniformly, the same move machinery and the same conservation
discipline apply without special cases.}

\clause{C-4.3}{Units divide into valued and non-valued. A unit is \textbf{valued}
if and only if it represents an economic claim whose value exists independently of
the governance or continued operation of the instrument that created it. A stock,
a bond, and a client's redemption claim on a managed account are valued. A
mandate's rulebook, a portfolio-swap reset mechanism, or a quantitative investment
strategy's governance shell is non-valued: its price is undefined --- not zero ---
because pricing the governance structure beside the claim it governs would count
the same economic reality twice.}

\clause{C-4.4}{Every unit carries three things: its terms (the contractual text,
versioned and append-only), its state (where it stands in its lifecycle), and its
smart contracts --- the executable form of those terms. From its terms follow its
\textbf{watches}: the dates, observations, and conditions the Event Monitor must
observe on its behalf. Declaring them is the unit's duty --- at registration, and
again whenever a later transaction creates a new obligation --- and the
declaration is exhaustive: the Monitor does exactly what the units have asked of
it, nothing more and nothing less.}

\clause{C-4.5}{A \textbf{wallet} is a named logical partition of position space ---
a container, and deliberately nothing more. Its state at any moment is a signed
quantity, its balance, for each unit it holds. A wallet names two parties: a
\textbf{manager}, who is authorised to perform transactions on it, and a
\textbf{beneficial owner}, who is legally accountable for the book. Wallets carry
no smart contract and no logic of their own; where a wallet appears to need
behaviour --- the rules of a managed account, for instance --- that behaviour is a
unit (the mandate).}

\clause{C-4.6}{An \textbf{atomic move} is the sole operation that changes a
balance: it transfers a positive quantity of exactly one unit from a source wallet
to a destination wallet in a single indivisible step. Quantities are exact
integers in minor units; where an event divides unevenly, the rounding convention
is part of the declared terms and the remainder is booked explicitly as its own
move --- nothing is ever silently dropped. Because a move only transfers, it can
never create or destroy quantity.}

\clause{C-4.7}{A \textbf{transaction} is an atomic, all-or-nothing bundle of four
kinds of change: moves, unit-state changes, position-state changes, and terms
changes. The four align one-for-one with the four components of the ledger's
state: balances, and the three homes of Section 7. A transaction applies in full
or not at all. Some transactions carry no moves (for example, the registration of
a new unit or the recording of a barrier breach).}

\clause{C-4.8}{The \textbf{immutable log} is the ordered, append-only,
hash-chained record of every admitted transaction --- hash-chained so that any
alteration of history is detectable. It is the sole canonical source of truth.
Three words are used precisely throughout this document: the \textbf{log} is the
recorded sequence itself; the \textbf{ledger} is the state obtained by folding the
log; and the \textbf{record} is everything that has been recorded --- the log and
all it carries, transactions, events, and observations alike. ``On the record''
means present there, and a projection is any view computed from it.}

\clause{C-4.9}{Every counterparty outside the system is represented by a virtual
wallet inside it, so the system is closed. Conservation holds by construction, not
by check.}

\clause{C-4.10}{When a balance must carry more information than a single number, it
generalises to a signed vector of coordinates --- owned, lent, posted --- whose
named rays are: lent out and borrowed, the two signs of lent; posted as collateral
and received as collateral, the two signs of posted.}

\pk{F10 (PARK-3), amends \S4}{%
  ``When a balance must carry more information than a single number, it generalises
  to a signed vector of coordinates --- owned, lent, posted --- whose named rays
  are: lent out and borrowed, the two signs of lent; posted as collateral and
  received as collateral, the two signs of posted.''}{%
  ``When a balance must carry more information than a single number, it generalises
  to a signed vector indexed by (coordinate, agreement): the coordinates owned,
  lent, posted, with named rays lent-out/borrowed and posted/received, each
  non-owned coordinate qualified by the collateral or margin agreement under which
  the quantity sits. Conservation holds per (unit, coordinate, agreement); the
  three-coordinate vector is the aggregate projection that sums over agreements.''}

\clause{C-4.11}{A quantity earns a coordinate only when a distinct real-world
action can change that coordinate independently of ownership. Only the owned
coordinate carries economic value and drives profit and loss. The remaining
coordinates carry mass, never value: they record possession and encumbrance under
a named agreement, no quantity reaches them without reference to the agreement
unit that governs it, and a lifecycle event extinguishes value, never mass --- a
unit leaves the ledger only from the zero vector. Anything computable from the
coordinates is a projection, computed when needed and never stored.}

% ===================================================================
\section*{\S5 --- The Machines}
% ===================================================================

\clause{C-5.1}{The framework has three machines --- one for each arrow of the
picture in Section 3. All three are generic machinery carrying no product
knowledge, and only the last ever writes the ledger: three roles, one writer.}

\clause{C-5.2}{The \textbf{Event Monitor} owns time and attention. It holds every
declared watch and evaluates them against the clock and the recorded observations,
emitting an event onto the record the moment a condition is met. It holds no
private list --- the watches it evaluates are on the record --- and it emits
events, never transactions. The relationship between the units and the Monitor is
a recorded handshake. Each declared watch is acknowledged by the Monitor, and the
acknowledgement --- an event like any other, flowing through the same pipeline ---
is written into the unit's state; a watch is thus visibly \emph{declared}, then
\emph{armed}, then \emph{fired} or \emph{expired}.}

\clause{C-5.3}{A declared watch that is never armed is therefore a visible gap,
never a silent one. The full chain is on the record and can be walked in either
direction: register a unit, the conditions it declares, the acknowledgements, the
events, the transactions, the state updates. The Monitor can affect timeliness and
liveness, but it cannot affect the correctness of ledger state, because it only
emits events --- it never proposes transactions.}

\clause{C-5.4}{The \textbf{Events Executor} owns delivery. It captures and records
the events arriving from outside, and, for each event --- arrived or emitted ---
it fires the responsible smart contract, handing it the event together with the
current projection of the ledger, and collects the proposed transaction. It runs
the contracts; it never writes the ledger; its output is only proposals.}

\clause{C-5.5}{The \textbf{Transaction Executor} owns the record. It is the
gatekeeper of the ledger door: the single component through which every
transaction must pass, and the only component permitted to change the ledger's
state. It validates structure --- conservation is automatic; authorisation,
idempotence, consistency of reference, and writer discipline are enforced --- and,
if the transaction is admissible, appends it to the immutable log. It creates
nothing and runs nothing; it admits or refuses. There is no other way to change
the ledger's state.}

\clause{C-5.6}{Every proposed transaction carries an identifier derived from its
cause, so a retried or duplicated proposal is recognised and committed only once.}

\clause{C-5.7}{The single-writer discipline is a correctness requirement. The
Transaction Executor therefore presents to the log a strictly monotonic sequence
of committed transactions; where two simultaneous events on the same instrument do
not commute and declare no precedence, it refuses rather than guesses an order.}

\clause{C-5.8}{The roles sit at different levels of trust. The Transaction Executor
is the trusted component: the ledger's integrity rests on it and on nothing else.
The contracts, the Event Monitor, and the Events Executor are outside the trust
boundary, untrusted in different ways: a contract can be economically wrong, so
its output is verified by recomputation; the Monitor and the Events Executor
cannot corrupt the ledger at all --- a late emission or firing produces the
identical transaction, a duplicate is committed once, a spurious one finds nothing
due and is refused --- but the timeliness of the system depends on them. Integrity
owes them nothing; liveness is conditional on them.}

% ===================================================================
\section*{\S6 --- Smart Contracts}
% ===================================================================

\clause{C-6.1}{Financial instruments are not passive data records; they are active
contracts that, when triggered by an event, emit the transactions reflecting the
economic consequences of that event. The Ledger performs none of this
product-specific logic: it admits well-formed transactions and records them. Since
the ledger is only a transaction processor, anything coming from outside must be
converted into a transaction, and the smart contract is the converter.}

\clause{C-6.2}{Everything a smart contract consumes must itself be on the record.}

\clause{C-6.3}{A contract never waits and never remembers: the watch waits in the
record, the Event Monitor notices when its condition is met, the Events Executor
delivers the firing --- and whatever a contract must carry from one firing to the
next, it records.}

\clause{C-6.4}{Every smart contract is a pure, deterministic function of its
inputs: the declared terms of the instrument, the recorded observations, and the
visible unit and position state at the moment of invocation. Each contract is
therefore a stateless, machine-readable encoding of the legal terms and conditions
of the instrument it governs.}

\clause{C-6.5}{The alignment with the legal contract runs deeper: every right and
obligation the legal contract creates is represented in the ledger as a unit. A
subscription is booked as an exchange --- cash against the client's redemption
claim --- never as a transfer for nothing.}

\clause{C-6.6}{Treating the client's redemption claim as a valued unit has direct
consequences for fee accrual and NAV attribution in managed accounts; these
consequences are resolved in the detailed specification consistently with the
principles stated here.}

% ===================================================================
\section*{\S7 --- State: The Three Homes}
% ===================================================================

\clause{C-7.1}{Every unit progresses through a well-defined sequence of lifecycle
states. Some elements of state are properties of the unit alone; others are joint
properties of the unit and a specific wallet.}

\clause{C-7.2}{The three-home state model encodes the distinction: (i) ProductTerms
--- the immutable, append-only versions of an instrument's contractual terms; (ii)
UnitStatus --- a materialised, replayable projection of events that affect the unit
as a whole; and (iii) PositionState --- per-wallet, per-unit attributes.}

\clause{C-7.3}{Because all three are derived from the immutable log rather than
being independent mutable stores, they can be discarded and rebuilt by replay at
any time without loss of information. Each non-balance fact has exactly one legal
writer.}

\clause{C-7.4}{The three homes are not bookkeeping conveniences: together with the
balances, they must carry \textbf{all the information the smart contracts require
in order to execute}. Each such fact is written into its home by an earlier
firing, so that the next firing finds its entire context on the record; if a
contract would ever need to remember something, that something has a home.}

\clause{C-7.5}{One event --- a settlement print, a margin payment, a split ---
lands in exactly one home, and the home is determined by what the fact is a
property of: the unit alone, the unit-and-holder pair, or the contract's terms.}

% ===================================================================
\section*{\S8 --- Valuation and Profit and Loss}
% ===================================================================

\clause{C-8.1}{The initial NAV of a wallet is the cash it took to build the
portfolio at inception; the current NAV is the cash you would hold if you unwound
every position at prevailing prices; and the profit and loss (PnL) is the
difference between the two.}

\clause{C-8.2}{At any point in time the NAV of a wallet is obtained by (i)
replaying the transaction log to recover the wallet's current composition and (ii)
multiplying that composition by the current price vector supplied by the pricing
data layer, the sum running over the valued units only.}

\clause{C-8.3}{Because unit state affects value, the pricing function is also
parameterised by the current lifecycle state of the instrument --- obtained from
the ledger, never from a store of its own.}

\clause{C-8.4}{A central and non-negotiable property follows: the profit and loss
of any wallet is exactly the variation in its NAV. PnL between two dates is
path-independent, and no second copy of ``realised PnL'' or ``book cost'' is ever
maintained alongside the position record.}

\clause{C-8.5}{For the variation in NAV to be exactly profit, deposits and
contributions must not create profit. Every inflow of value is one of three
things:}

\clause{C-8.6}{an exchange paired with an equal-valued outflow --- settled-to-market
margin is this case, a settlement extinguishing the day's exposure with no return
obligation;}

\clause{C-8.7}{a contribution or financing received against an equal obligation
created in the same transaction --- collateral received under title transfer, cash
included, is this case, owned by the receiver against an equivalent-return
obligation that is itself a unit;}

\pk{F12 (PARK-4), amends \S8 case 2}{%
  ``a contribution or financing received against an equal obligation created in the
  same transaction --- collateral received under title transfer, cash included, is
  this case, owned by the receiver against an equivalent-return obligation that is
  itself a unit''.}{%
  ``a contribution or financing received against an obligation created in the same
  transaction and valued at fair value, any day-one difference between the inflow
  amount and the obligation's fair value being recognised as financing basis ---
  collateral received under title transfer, cash included, is this case, owned by
  the receiver against an equivalent-return obligation that is itself a unit''.}

\clause{C-8.8}{or value held in custody without being owned, recorded on the
collateral-received coordinate --- collateral received under a security interest
without right of use, segregated cash included, is this case.}

\clause{C-8.9}{Where an inflow moves under a margin or collateral agreement, which
case governs it is declared once, on that agreement unit. A deposit therefore
cannot move profit and loss --- not because a formula subtracts it, but because it
cannot change net owned value.}

% ===================================================================
\section*{\S9 --- Corporate Actions and the Market Data Operator}
% ===================================================================

\clause{C-9.1}{A corporate action is, first of all, a transaction like any other: a
list of moves, together with unit-state changes and a new version of the
instrument's terms, all applied atomically through the same Transaction Executor.}

\clause{C-9.2}{A corporate action also changes the meaning of all the market data
recorded before it. To carry pre-event market data into the post-event frame, a
transformation must be applied. The framework gives it a name: the \textbf{market
data operator}. Every kind of market data is equipped with its own operators, one
per kind of event --- declared once, when the kind of data is introduced and the
corporate action is recorded, never improvised at the moment of reading.}

\clause{C-9.3}{Three principles govern the operators. First, the operator is
intrinsic to the pairing of the kind of data and the kind of event. Second, the
ledger is the authority; the pricing data layer applies. Third, originals are
never overwritten. The observed price stays on the record exactly as observed; the
adjusted value is computed at each read. Derived quantities are recomputed from
operator-adjusted inputs.}

% ===================================================================
\section*{\S10 --- Strategies, Virtual Ledgers, and Total Return Swaps}
% ===================================================================

\clause{C-10.1}{A \textbf{virtual ledger} is a complete instance of the same
machine --- units, wallets, balances, an immutable log, the three machines ---
whose holdings are notional: it settles nothing, and none of its wallets ever
reaches a settlement instruction or the balance sheet. Virtual ledgers are subject
to the same six commitments as true ledgers, especially full reproducibility and
time travel. The word virtual keeps one meaning: on the far side of the settlement
boundary.}

\clause{C-10.2}{A deterministic strategy is a smart contract whose rulebook is the
terms of a (non-valued) strategy unit. The hypothetical portfolio it manages is a
wallet in a virtual ledger. Each rebalancing is a recorded transaction triggered
by stamped market observations. The strategy's level is that wallet's NAV --- a
projection computed by the same fold as any other; because the strategy ledger is
replayable, any party holding the recorded rules and observations recomputes the
level exactly, and an index restatement is a compensating entry, never a rewrite.}

\clause{C-10.3}{Two disciplines make this safe. First, the bridge between ledgers
is an observation, never a move. Second, the wall between models and the ledger
holds level by level: the true ledger never runs the strategy; it consumes the
level as an observation.}

\clause{C-10.4}{When the strategy's trades are real, the strategy's output toward
the world is an order (an instruction, not a fact). The executions that come back
are external events, mapped by contracts and folded into a real wallet of the true
ledger. The virtual ledger becomes the benchmark: it holds what the rules say the
portfolio should be, the real wallet holds what execution achieved, and the
difference between the two folds is the slippage --- exact, never estimated.}

\clause{C-10.5}{A total return swap is then one definition: the reference leg of a
TRS is the NAV of a designated wallet (or virtual ledger). In every case the TRS
is an ordinary unit in the true ledger whose smart contract, at each reset, reads
the stamped NAV observation and emits the reset and financing moves.}

% ===================================================================
\section*{\S11 --- Structural Invariants and the Impossibility of Broken States}
% ===================================================================

\clause{C-11.1}{The Ledger is deliberately engineered so that certain classes of
inconsistent or illegal state cannot be represented at all.}

\clause{C-11.2}{\textbf{Atomicity.} A transaction is applied in full or not at
all.}

\clause{C-11.3}{\textbf{Consistency of reference.} A quantity and a price may be
combined only when they refer to the same state of the instrument. After a
two-for-one split, 1,000 shares at 100 become 2,000 at 50; multiplying the new
count of 2,000 by the stale price of 100 would book a phantom valuation.}

\clause{C-11.4}{\textbf{Writer discipline.} Lifecycle facts are written only by the
single legal writer for that fact.}

\clause{C-11.5}{These three join the invariants established earlier ---
conservation by construction (Section 4), the append-only log (Section 4), the
single writer (Section 5), the determinism of every arrow (Section 3), and
idempotence under the cause-derived identifier (Sections 5 and 12) --- to form the
catalogue that the property-based tests of the Testability commitment exercise.
Together they render the classes of internal reconciliation failure that dominate
conventional architectures unreachable within the Ledger's scope: what cannot be
represented cannot break.}

% ===================================================================
\section*{\S12 --- Time Travel, Replayability, and Idempotence}
% ===================================================================

\clause{C-12.1}{Every historical state of the Ledger is exactly reconstructible by
replaying the immutable transaction log up to the desired timestamp.}

\pk{F2 (PARK-2), amends \S12}{%
  ``Every historical state of the Ledger is exactly reconstructible by replaying
  the immutable transaction log up to the desired timestamp.''}{%
  ``Every historical state of the Ledger is exactly reconstructible from the
  immutable transaction log along two axes: a transaction-time horizon (the log
  prefix --- what was known) and a valid-time horizon (which recorded observations
  are in force). Fixing both yields a deterministic projection; the as-known-at
  view over a prefix is never rewritten by a later correction, and the as-of view
  recomputes bit-exactly from the prefix that includes the correction.''}

\clause{C-12.2}{Because it reads only committed transactions and never re-executes
smart contracts, the historical record is independent of the machinery that
produced it.}

\clause{C-12.3}{Smart contracts must be idempotent. The mechanism is the
cause-derived identifier of Section 5: a retried or late-arriving proposal for an
already-processed event produces no duplicate effect.}

% ===================================================================
\section*{\S13 --- Two Layers of Correctness}
% ===================================================================

\clause{C-13.1}{\textbf{Structural correctness} is guaranteed by the Transaction
Executor. These properties hold regardless of whether the smart contract that
produced the transaction was correct, careless, or malicious.}

\clause{C-13.2}{\textbf{Economic correctness} is not guaranteed at admission; it is
made checkable by recomputation. A smart contract may emit a perfectly well-formed
transaction that nevertheless books the wrong economics. The discrepancy is
detected after the fact by recomputing the transaction from the recorded inputs.
The difference constitutes a defect of the contract, not of the ledger, and it is
repaired by a subsequent compensating transaction through the same door.}

\clause{C-13.3}{The Ledger relies on structural correctness because it is
guaranteed by construction; it expects economic correctness of its contracts and
renders it mechanically verifiable. Detection is after the fact. The framework
prefers a detectable, repairable error over a pretence of prevention it cannot
deliver.}

% ===================================================================
\section*{\S14 --- Minimum Requirements}
% ===================================================================

\subsection*{\S14.1 The Event Monitor and the Events Executor}

\clause{C-14.1}{The Event Monitor and the Events Executor are infrastructure
outside the ledger's trust boundary. Any implementation must satisfy, at
minimum:}

\clause{C-14.2}{Durable watches and timers. A watch registered for a future date or
condition fires when it is met, regardless of restarts or failures; timer state is
persisted.}

\clause{C-14.3}{At-least-once emission. No event is ever lost; a duplicate is
harmless because of idempotence checking.}

\clause{C-14.4}{Acknowledged watches. Every declared watch is acknowledged, and the
acknowledgement is recorded in the unit's state; a watch that remains
unacknowledged beyond its arming window is a visible, overdue item, never
silence.}

\clause{C-14.5}{Triggers carry timing, not data. A trigger says when; everything
the event means is read from the record.}

\clause{C-14.6}{Deterministic, replayable orchestration. Each machine's own state
is reconstructible by replaying its append-only history.}

\clause{C-14.7}{No privilege. Neither machine holds write access to the ledger.}

\clause{C-14.8}{Obligation liveness. Every contractual obligation is registered
with a deadline, a test for its discharge, and a fallback action, and by its
deadline it is discharged, compensated, or declared defaulted; a duty that fails to
fire is a visible, overdue item, never silence.}

\subsection*{\S14.2 The Pricing Stack and the Pricing Data Layer}

\clause{C-14.9}{The ledger records the outputs of pricing models and never runs
one.}

\clause{C-14.10}{Valuation is state-dependent and obtains that state from the
ledger.}

\clause{C-14.11}{Valuations are reproducible from recorded inputs.}

\clause{C-14.12}{Prices and quantities are never mixed across an event.}

\clause{C-14.13}{Originals are preserved; adjusted values are computed; derived
values are recomputed.}

\clause{C-14.14}{Estimates and entitlements are kept apart.}

\clause{C-14.15}{Model outputs re-enter only as new observations, carrying
sufficient provenance. Reproducing a valuation from recorded inputs and identified
prices is unconditional; reproducing a model's number requires the model itself,
whose retention is a governance matter outside the ledger's scope.}

% ===================================================================
\section*{Scope of the Ledger Project}
% ===================================================================

\clause{C-Scope.1}{The Ledger project encompasses the design, formal
specification, reference implementation, and verification of the following
components:}

\clause{C-Scope.2}{The closed double-entry move model, the conservation law, and
the single-door admission function.}

\clause{C-Scope.3}{The immutable, hash-chained transaction log and the projection
machinery that derives all views from it.}

\clause{C-Scope.4}{The three-home state architecture (ProductTerms, UnitStatus,
PositionState) and the conditions that render illegal states unrepresentable.}

\clause{C-Scope.5}{The smart-contract framework, its event-triggering discipline,
and the purity and idempotence requirements.}

\clause{C-Scope.6}{The market data operators and the integration layer between
lifecycle state, pricing models, and the corporate-action adjustment machinery.}

\clause{C-Scope.7}{The time-travel and historical-replay machinery, with its
guarantees of determinism and reproducibility.}

\clause{C-Scope.8}{The formal invariant catalogue and the property-based testing
regime.}

\clause{C-Scope.9}{The settlement-projection layer.}

\clause{C-Scope.10}{The virtual-ledger mechanism: subordinate, non-settling
instances of the same machine whose projections enter the true ledger only as
stamped observations.}

\clause{C-Scope.11}{Explicitly out of scope are: price discovery and executable
pricing; legal enforceability or validity of contracts; settlement finality,
payment-system access, and CSD connectivity; regulatory-submission gateways and
report formatting; model governance and validation; and accounting-policy
decisions. These concerns lie outside the Ledger boundary and are reconciled
against it only at its perimeter.}

% ===================================================================
\section*{Authority and Amendment}
% ===================================================================

\clause{C-Auth.1}{This manifesto is the authoritative statement of intent and scope
for the Ledger project --- its architectural constitution --- and the direction of
authority is one-way: the manifesto is rated against nothing; everything else is
rated against it.}

\clause{C-Auth.2}{All specifications, implementations, tests, and integration
layers --- including every existing version of the ledger specification --- shall
be assessed for conformance against it and brought into line where they deviate.}

\clause{C-Auth.3}{Where a genuine improvement requires departing from this
document, the manifesto is amended explicitly first; a design choice that cannot be
traced to its principles shall be rejected.}

\clause{C-Auth.4}{Its vocabulary --- unit, wallet, balance, move, transaction,
watch, the immutable log, projection, the Event Monitor, the Events Executor, the
Transaction Executor, smart contract, the market data operator, the three homes,
virtual wallet, virtual ledger --- is fixed: one name per component, and no
specification shall introduce a synonym for any of them.}

\vfill
\begin{center}\small\itshape
--- End of clause-addressed rendering. 110 normative clauses. Three parked
amendments (F2, F10, F12) shown as marked v1.2 deltas. Not one word of
Constitution v1.1 prose changed. ---
\end{center}

\end{document}
